\section{Discussion}

\subsection{Systematic Comparison with State-of-the-Art Estimators}
Table~\ref{tab:sota_comparison} compares the proposed framework against classical non-parametric and modern deep learning methodologies across six core dimensions.

\begin{table*}[!t]
\caption{Systematic Comparison with State-of-the-Art Transfer Entropy Estimators}
\label{tab:sota_comparison}
\centering
\begin{tabular}{lcccc}
\toprule
Evaluation Dimension & Classical KSG \cite{kraskov2004estimating} & TREET (2024/2025) \cite{treet2024} & TENDE (2026) \cite{tende2025} & \textbf{Proposed ANTE-UDA (This Work)} \\
\midrule
Target Observation Window & Large ($N > 100$) & Long ($T \ge 1,000$) & Massive ($T \ge 50,000$) & \textbf{Ultra-short ($N = 10 - 100$)} \\
Computational Complexity & $\mathcal{O}(k \cdot N \log N)$ & $\mathcal{O}(T^2)$ Attention & Diffusion ($K \approx 100$ steps) & \textbf{$\mathcal{O}(1)$ Forward Pass (7.0 ms)} \\
Finite-Sample Variance & High ($k$-NN dispersion) & Explodes on short series & Unstable on short series & \textbf{Provably Reduced (Prop.~\ref{prop:var_reduction})} \\
Domain Shift Adaptation & None (Non-parametric) & Supervised re-training & Score matching & \textbf{Label-Free Variational UDA} \\
Noise Robustness (AWGN) & Metric distortion & Unverified & Implicit diffusion & \textbf{Smooth Activation Regularization} \\
Deployment Viability & CPU-bound (Slow) & GPU Server Required & GPU Cluster Required & \textbf{Edge-AI \& Microcontrollers} \\
\bottomrule
\end{tabular}
\end{table*}

\subsection{Why Transformers and Diffusion Models Falter on Short Windows}
Recent advancements in information theory have produced powerful neural estimators, yet their design paradigms conflict with physiological realities:
\begin{enumerate}
    \item \textbf{Transformer Attention Dispersion in TREET \cite{treet2024}:} TREET optimizes the Donsker-Varadhan bound over causal attention weights. While effective when historical context is long ($l \ge 99$), its multi-head attention entropy collapses when applied to ultra-short windows ($N \le 50$). Because the sequence length is insufficient to populate the softmax attention matrix, gradient estimates become degenerate, leading to severe negative estimation collapse (e.g., $-3.02 \pm 2.59$ on uncoupled synthetic benchmarks).
    \item \textbf{Computational Intractability of Diffusion in TENDE \cite{tende2025}:} TENDE successfully bypasses the McAllester-Stratos exponential sample size lower bound $\mathcal{O}(e^I)$ \cite{mcallester2020formal} by converting KL divergence into a continuous score-matching integral over a Variance-Preserving SDE. However, evaluating this integral requires numerically solving the reverse-time SDE across hundreds of discretization steps ($K \approx 100$). Benchmarks indicate an evaluation latency of \textbf{approximately 20 minutes per time-series pair on an NVIDIA A100 GPU}. For wearable continuous health monitoring, this computational overhead is prohibitive.
    \item \textbf{The Edge Advantage of ANTE-UDA:} By restricting model capacity to a compact 369-parameter statistics network, our amortized approach completes inference in \textbf{7.0~ms per window on standard CPUs}, achieves provable variance reduction via parameter sharing, and adapts to patient physiology without ground-truth labels.
\end{enumerate}

\subsection{Analytical Insights into Hybrid Quantum Circuits}
Decomposing estimator error into $\mathrm{MSE} = \mathrm{Bias}^2 + \mathrm{Variance}$ reveals profound behavioral contrasts between classical and quantum representations:
\begin{itemize}
    \item \textbf{Exceptional Quantum Variance Suppression:} The Hybrid Classical-Quantum circuit (MLP + 4-qubit VQC) achieves an estimation variance of \textbf{$1.77 \times 10^{-5}$}, which is \textbf{160 times lower than classical MLP} ($2.88 \times 10^{-3}$). Unitary rotations and circular CNOT entanglement map classical coordinates into an orthonormal Hilbert space $\mathbb{C}^{16}$, providing superior geometric rigidity against stochastic noise.
    \item \textbf{Simulation-Induced Underfitting Bias:} The higher total MSE of Hybrid VQC ($0.0274$) is dominated by estimation bias ($\text{Bias} \approx -0.165$). This bias is an artifact of the 7-epoch training budget imposed by classical CPU statevector simulation of quantum circuits.
    \item \textbf{Strategic Architecture Positioning:} In clinical deployment, the 369-parameter Small MLP serves as the primary real-time engine. Meanwhile, the Hybrid VQC establishes a foundational proof-of-concept for quantum causal discovery, demonstrating that quantum embedding offers unprecedented variance suppression for noisy biological processes.
\end{itemize}

\subsection{Eliminating Algorithmic Selection Artifacts in Aging Research}
A critical finding of our clinical investigation is the resolution of a longstanding algorithmic selection bias in autonomic physiology literature. Previous studies enforced a heuristic synchronization threshold ($TE \ge 0.02$ nats) to filter out ``poorly coupled'' segments, inadvertently discarding 17 out of 20 elderly subjects (85\% of the Old cohort). Our full-cohort evaluation confirms that diminished Transfer Entropy in elderly populations is not a signal acquisition flaw, but rather the genuine physiological manifestation of autonomic vagal attenuation and age-related cardiorespiratory decoupling.

\subsection{Clinical Feasibility and Wearable Translation}
The sub-10~ms inference latency and minimal memory footprint ($\approx 1.5$~KB) of our Small MLP make it suitable for on-chip integration into low-power biomedical wearables (e.g., smartwatches, Holter monitors). By continuously tracking directed information flow $\Delta TE$, wearable devices can deliver real-time alerts for autonomic decompensation, sleep apnea episodes, and cardiovascular distress.
